% --- Archon physics formalization source begin ---
% archon:physics
% archon:covers PhyXMiniProblems/problem_phyx_mini_0307.lean
% archon:source-report reports/phyx_mini/problem_phyx_mini_0307.source.json

\chapter{Physics problem phyx\_mini\_0307}
\label{ch:PhyXMiniProblems_problem_phyx_mini_0307}

\paragraph{Problem source.}
\#\# Physical scenario

A continuous traveling wave with amplitude \$A\$ is incident on a boundary. The continuous reflection, with a smaller amplitude \$B\$, travels back through the incoming wave. The resulting interference pattern is displayed in figure. The standing wave ratio is defined to be \$\$SWR = \textbackslash{}frac\{A + B\}\{A - B\}.\$\$ The reflection coefficient \$R\$ is the ratio of the power of the reflected wave to the power of the incoming wave and is thus proportional to the ratio \$(B/A)\textasciicircum{}2\$.

\#\# Question

For SWR = \$1.50\$, what is \$R\$ expressed as a percentage?

\#\# Answer choices

- A: A: 3.2\%

- B: B: 3.5\%

- C: C: 3.8\%

- D: D: 4.0\%

\#\# Figure caption (auxiliary; use the image as primary evidence)

The image shows a wavy green line resembling a sinusoidal wave with peaks and troughs. Three vertical, black double-headed arrows mark specific points on the wave:

1. The first arrow points between a peak and a trough, labeled "A\_max" (maximum amplitude).
2. The middle arrow points between a trough and a peak, labeled "A\_min" (minimum amplitude).
3. The third arrow, similar to the first, points between a peak and a trough and is also labeled "A\_max."

The green wave appears to smoothly oscillate, illustrating amplitude changes across its length.

\#\# Dataset metadata

Wave Properties; Multi-Formula Reasoning, Numerical Reasoning

\paragraph{Recorded answer/context.}
D: D: 4.0\%

\paragraph{Figure/image path.}
/root/proposal\_for\_physic/science-mango/phyx\_mini\_run/phyx\_data/test\_image/307.png

\paragraph{Formalization target.}
create a compiling Lean file with sorry bodies at `PhyXMiniProblems/problem\_phyx\_mini\_0307.lean`. The Lean declarations must preserve the physical quantities, dimensions or dimensional roles, figure labels, governing-law hypotheses, and final relation expressed by this problem.
Use Mathlib/Physlib names found through LeanExplore where available. If a physics API is missing, introduce faithful local abstractions rather than scalar placeholder aliases.

\begin{theorem}[Physics formalization target]
\label{thm:physics:phyx_mini_0307:target}
\lean{PhyXMiniProblems.ProblemPhyXMini0307.problem_phyx_mini_0307}
\uses{def:physics:phyx-mini-0307:phyxminiproblems-problemphyxmini0307-boundaryreflectionsetup, def:physics:phyx-mini-0307:phyxminiproblems-problemphyxmini0307-matchesproblemandfiguredata, def:physics:phyx-mini-0307:phyxminiproblems-problemphyxmini0307-hasphysicalwaveparameters, def:physics:phyx-mini-0307:phyxminiproblems-problemphyxmini0307-satisfiesstandingwaveinterferencelaw, def:physics:phyx-mini-0307:phyxminiproblems-problemphyxmini0307-satisfiescommonmediumpowerreflectionlaw, def:physics:phyx-mini-0307:phyxminiproblems-problemphyxmini0307-reflectioncoefficientpercentage, def:physics:phyx-mini-0307:phyxminiproblems-problemphyxmini0307-matchesanswerchoice, def:physics:phyx-mini-0307:phyxminiproblems-problemphyxmini0307-recordeddatasetanswer}
The assigned autoformalize agent should translate this physics problem into Lean declarations in the covered file, with theorem and lemma proof bodies written as `by sorry`.
\end{theorem}
\begin{proof}
This is an autoformalization task, not a proof task. Produce statements that can later be proved by the physics prover without weakening the physical model.
\end{proof}

% --- Archon declaration topology begin ---
\section{Declaration topology}
\begin{definition}[power Dimension]
  \label{def:physics:phyx-mini-0307:phyxminiproblems-problemphyxmini0307-powerdimension}
  \lean{PhyXMiniProblems.ProblemPhyXMini0307.powerDimension}
  Physical dimension of power: mass times length squared per time cubed
\end{definition}
\begin{proof}
  This is the declared model definition of power Dimension.
\end{proof}
\begin{definition}[Wave Amplitude]
  \label{def:physics:phyx-mini-0307:phyxminiproblems-problemphyxmini0307-waveamplitude}
  \lean{PhyXMiniProblems.ProblemPhyXMini0307.WaveAmplitude}
  A nonnegative, unit-independent amplitude of a wave of unspecified kind
\end{definition}
\begin{proof}
  This is the declared model definition of Wave Amplitude.
\end{proof}
\begin{definition}[Wave Power]
  \label{def:physics:phyx-mini-0307:phyxminiproblems-problemphyxmini0307-wavepower}
  \lean{PhyXMiniProblems.ProblemPhyXMini0307.WavePower}
  \uses{def:physics:phyx-mini-0307:phyxminiproblems-problemphyxmini0307-powerdimension}
  A nonnegative, unit-independent wave power
\end{definition}
\begin{proof}
  This is the declared model definition of Wave Power.
\end{proof}
\begin{definition}[Wave Speed]
  \label{def:physics:phyx-mini-0307:phyxminiproblems-problemphyxmini0307-wavespeed}
  \lean{PhyXMiniProblems.ProblemPhyXMini0307.WaveSpeed}
  A nonnegative, unit-independent propagation speed
\end{definition}
\begin{proof}
  This is the declared model definition of Wave Speed.
\end{proof}
\begin{definition}[amplitude SIReadout]
  \label{def:physics:phyx-mini-0307:phyxminiproblems-problemphyxmini0307-amplitudesireadout}
  \lean{PhyXMiniProblems.ProblemPhyXMini0307.amplitudeSIReadout}
  \uses{def:physics:phyx-mini-0307:phyxminiproblems-problemphyxmini0307-waveamplitude}
  Coherent SI scalar readout of an amplitude
\end{definition}
\begin{proof}
  This is the declared model definition of amplitude SIReadout.
\end{proof}
\begin{definition}[power In Watts]
  \label{def:physics:phyx-mini-0307:phyxminiproblems-problemphyxmini0307-powerinwatts}
  \lean{PhyXMiniProblems.ProblemPhyXMini0307.powerInWatts}
  \uses{def:physics:phyx-mini-0307:phyxminiproblems-problemphyxmini0307-wavepower}
  Watt readout of a wave power
\end{definition}
\begin{proof}
  This is the declared model definition of power In Watts.
\end{proof}
\begin{definition}[speed In Meters Per Second]
  \label{def:physics:phyx-mini-0307:phyxminiproblems-problemphyxmini0307-speedinmeterspersecond}
  \lean{PhyXMiniProblems.ProblemPhyXMini0307.speedInMetersPerSecond}
  \uses{def:physics:phyx-mini-0307:phyxminiproblems-problemphyxmini0307-wavespeed}
  Metres-per-second readout of a propagation speed
\end{definition}
\begin{proof}
  This is the declared model definition of speed In Meters Per Second.
\end{proof}
\begin{definition}[Horizontal Direction]
  \label{def:physics:phyx-mini-0307:phyxminiproblems-problemphyxmini0307-horizontaldirection}
  \lean{PhyXMiniProblems.ProblemPhyXMini0307.HorizontalDirection}
  Horizontal directions used only to distinguish the two propagation senses
\end{definition}
\begin{proof}
  This is the declared model definition of Horizontal Direction.
\end{proof}
\begin{definition}[opposite Direction]
  \label{def:physics:phyx-mini-0307:phyxminiproblems-problemphyxmini0307-oppositedirection}
  \lean{PhyXMiniProblems.ProblemPhyXMini0307.oppositeDirection}
  \uses{def:physics:phyx-mini-0307:phyxminiproblems-problemphyxmini0307-horizontaldirection}
  The direction opposite to a given horizontal propagation direction
\end{definition}
\begin{proof}
  This is the declared model definition of opposite Direction.
\end{proof}
\begin{definition}[Continuous Traveling Wave]
  \label{def:physics:phyx-mini-0307:phyxminiproblems-problemphyxmini0307-continuoustravelingwave}
  \lean{PhyXMiniProblems.ProblemPhyXMini0307.ContinuousTravelingWave}
  \uses{def:physics:phyx-mini-0307:phyxminiproblems-problemphyxmini0307-waveamplitude, def:physics:phyx-mini-0307:phyxminiproblems-problemphyxmini0307-wavepower, def:physics:phyx-mini-0307:phyxminiproblems-problemphyxmini0307-wavespeed, def:physics:phyx-mini-0307:phyxminiproblems-problemphyxmini0307-horizontaldirection}
  A traveling-wave component. The dimensionless profile gives its shape as a function of phase; continuity records the source's description of both the incident and reflected waves as continuous. Amplitude and power remain genuine physical quantities rather than scalar aliases
\end{definition}
\begin{proof}
  This is the declared model definition of Continuous Traveling Wave.
\end{proof}
\begin{definition}[Figure Arrow]
  \label{def:physics:phyx-mini-0307:phyxminiproblems-problemphyxmini0307-figurearrow}
  \lean{PhyXMiniProblems.ProblemPhyXMini0307.FigureArrow}
  The three vertical amplitude arrows visible in the primary image
\end{definition}
\begin{proof}
  This is the declared model definition of Figure Arrow.
\end{proof}
\begin{definition}[Envelope Label]
  \label{def:physics:phyx-mini-0307:phyxminiproblems-problemphyxmini0307-envelopelabel}
  \lean{PhyXMiniProblems.ProblemPhyXMini0307.EnvelopeLabel}
  Text labels printed beside the figure's amplitude arrows
\end{definition}
\begin{proof}
  This is the declared model definition of Envelope Label.
\end{proof}
\begin{definition}[Standing Wave Envelope Figure]
  \label{def:physics:phyx-mini-0307:phyxminiproblems-problemphyxmini0307-standingwaveenvelopefigure}
  \lean{PhyXMiniProblems.ProblemPhyXMini0307.StandingWaveEnvelopeFigure}
  \uses{def:physics:phyx-mini-0307:phyxminiproblems-problemphyxmini0307-waveamplitude, def:physics:phyx-mini-0307:phyxminiproblems-problemphyxmini0307-figurearrow, def:physics:phyx-mini-0307:phyxminiproblems-problemphyxmini0307-envelopelabel}
  The standing-wave envelope read from the bitmap. Arrow amplitudes are physical quantities; no pixel length is interpreted as a calibrated physical measurement
\end{definition}
\begin{proof}
  This is the declared model definition of Standing Wave Envelope Figure.
\end{proof}
\begin{definition}[Boundary Reflection Setup]
  \label{def:physics:phyx-mini-0307:phyxminiproblems-problemphyxmini0307-boundaryreflectionsetup}
  \lean{PhyXMiniProblems.ProblemPhyXMini0307.BoundaryReflectionSetup}
  \uses{def:physics:phyx-mini-0307:phyxminiproblems-problemphyxmini0307-continuoustravelingwave, def:physics:phyx-mini-0307:phyxminiproblems-problemphyxmini0307-standingwaveenvelopefigure}
  The boundary-reflection experiment. The dimensionless standing-wave ratio and reflection coefficient are independent measured/model quantities here; neither is defined from the requested numerical answer
\end{definition}
\begin{proof}
  This is the declared model definition of Boundary Reflection Setup.
\end{proof}
\begin{definition}[incident Amplitude SI]
  \label{def:physics:phyx-mini-0307:phyxminiproblems-problemphyxmini0307-incidentamplitudesi}
  \lean{PhyXMiniProblems.ProblemPhyXMini0307.incidentAmplitudeSI}
  \uses{def:physics:phyx-mini-0307:phyxminiproblems-problemphyxmini0307-amplitudesireadout, def:physics:phyx-mini-0307:phyxminiproblems-problemphyxmini0307-boundaryreflectionsetup}
  SI readout of the incident amplitude A
\end{definition}
\begin{proof}
  This is the declared model definition of incident Amplitude SI.
\end{proof}
\begin{definition}[reflected Amplitude SI]
  \label{def:physics:phyx-mini-0307:phyxminiproblems-problemphyxmini0307-reflectedamplitudesi}
  \lean{PhyXMiniProblems.ProblemPhyXMini0307.reflectedAmplitudeSI}
  \uses{def:physics:phyx-mini-0307:phyxminiproblems-problemphyxmini0307-amplitudesireadout, def:physics:phyx-mini-0307:phyxminiproblems-problemphyxmini0307-boundaryreflectionsetup}
  SI readout of the reflected amplitude B
\end{definition}
\begin{proof}
  This is the declared model definition of reflected Amplitude SI.
\end{proof}
\begin{definition}[maximum Envelope Amplitude SI]
  \label{def:physics:phyx-mini-0307:phyxminiproblems-problemphyxmini0307-maximumenvelopeamplitudesi}
  \lean{PhyXMiniProblems.ProblemPhyXMini0307.maximumEnvelopeAmplitudeSI}
  \uses{def:physics:phyx-mini-0307:phyxminiproblems-problemphyxmini0307-amplitudesireadout, def:physics:phyx-mini-0307:phyxminiproblems-problemphyxmini0307-boundaryreflectionsetup}
  Figure readout denoted by the left-hand A\_max label
\end{definition}
\begin{proof}
  This is the declared model definition of maximum Envelope Amplitude SI.
\end{proof}
\begin{definition}[minimum Envelope Amplitude SI]
  \label{def:physics:phyx-mini-0307:phyxminiproblems-problemphyxmini0307-minimumenvelopeamplitudesi}
  \lean{PhyXMiniProblems.ProblemPhyXMini0307.minimumEnvelopeAmplitudeSI}
  \uses{def:physics:phyx-mini-0307:phyxminiproblems-problemphyxmini0307-amplitudesireadout, def:physics:phyx-mini-0307:phyxminiproblems-problemphyxmini0307-boundaryreflectionsetup}
  Figure readout denoted by the central A\_min label
\end{definition}
\begin{proof}
  This is the declared model definition of minimum Envelope Amplitude SI.
\end{proof}
\begin{definition}[Matches Problem And Figure Data]
  \label{def:physics:phyx-mini-0307:phyxminiproblems-problemphyxmini0307-matchesproblemandfiguredata}
  \lean{PhyXMiniProblems.ProblemPhyXMini0307.MatchesProblemAndFigureData}
  \uses{def:physics:phyx-mini-0307:phyxminiproblems-problemphyxmini0307-amplitudesireadout, def:physics:phyx-mini-0307:phyxminiproblems-problemphyxmini0307-oppositedirection, def:physics:phyx-mini-0307:phyxminiproblems-problemphyxmini0307-boundaryreflectionsetup, def:physics:phyx-mini-0307:phyxminiproblems-problemphyxmini0307-maximumenvelopeamplitudesi}
  Direct prose and image evidence. The two component waves propagate in opposite directions, the bitmap labels its two outer arrows A\_max and its central arrow A\_min, and the supplied SWR readout is exactly 1.50 = 3/2. No reflection coefficient or answer percentage is asserted here
\end{definition}
\begin{proof}
  This is the declared model definition of Matches Problem And Figure Data.
\end{proof}
\begin{definition}[Has Physical Wave Parameters]
  \label{def:physics:phyx-mini-0307:phyxminiproblems-problemphyxmini0307-hasphysicalwaveparameters}
  \lean{PhyXMiniProblems.ProblemPhyXMini0307.HasPhysicalWaveParameters}
  \uses{def:physics:phyx-mini-0307:phyxminiproblems-problemphyxmini0307-powerinwatts, def:physics:phyx-mini-0307:phyxminiproblems-problemphyxmini0307-speedinmeterspersecond, def:physics:phyx-mini-0307:phyxminiproblems-problemphyxmini0307-boundaryreflectionsetup, def:physics:phyx-mini-0307:phyxminiproblems-problemphyxmini0307-incidentamplitudesi, def:physics:phyx-mini-0307:phyxminiproblems-problemphyxmini0307-reflectedamplitudesi}
  Nondegeneracy conditions from the physical scenario. In particular, the reflected amplitude B is smaller than the positive incident amplitude A
\end{definition}
\begin{proof}
  This is the declared model definition of Has Physical Wave Parameters.
\end{proof}
\begin{definition}[Satisfies Standing Wave Interference Law]
  \label{def:physics:phyx-mini-0307:phyxminiproblems-problemphyxmini0307-satisfiesstandingwaveinterferencelaw}
  \lean{PhyXMiniProblems.ProblemPhyXMini0307.SatisfiesStandingWaveInterferenceLaw}
  \uses{def:physics:phyx-mini-0307:phyxminiproblems-problemphyxmini0307-boundaryreflectionsetup, def:physics:phyx-mini-0307:phyxminiproblems-problemphyxmini0307-incidentamplitudesi, def:physics:phyx-mini-0307:phyxminiproblems-problemphyxmini0307-reflectedamplitudesi, def:physics:phyx-mini-0307:phyxminiproblems-problemphyxmini0307-maximumenvelopeamplitudesi, def:physics:phyx-mini-0307:phyxminiproblems-problemphyxmini0307-minimumenvelopeamplitudesi}
  Constructive and destructive interference give envelope amplitudes A + B and A - B. Their ratio is the standing-wave ratio. These are symbolic governing laws and contain no substituted SWR or requested percentage
\end{definition}
\begin{proof}
  This is the declared model definition of Satisfies Standing Wave Interference Law.
\end{proof}
\begin{definition}[Satisfies Common Medium Power Reflection Law]
  \label{def:physics:phyx-mini-0307:phyxminiproblems-problemphyxmini0307-satisfiescommonmediumpowerreflectionlaw}
  \lean{PhyXMiniProblems.ProblemPhyXMini0307.SatisfiesCommonMediumPowerReflectionLaw}
  \uses{def:physics:phyx-mini-0307:phyxminiproblems-problemphyxmini0307-powerinwatts, def:physics:phyx-mini-0307:phyxminiproblems-problemphyxmini0307-boundaryreflectionsetup, def:physics:phyx-mini-0307:phyxminiproblems-problemphyxmini0307-incidentamplitudesi, def:physics:phyx-mini-0307:phyxminiproblems-problemphyxmini0307-reflectedamplitudesi}
  Incident and reflected waves in the same medium share a power scale k, so each power is k times amplitude squared. The reflection coefficient is independently constrained to be reflected power divided by incident power. This exposes the proportionality invoked by the source without assuming its numerical consequence
\end{definition}
\begin{proof}
  This is the declared model definition of Satisfies Common Medium Power Reflection Law.
\end{proof}
\begin{definition}[Answer Choice]
  \label{def:physics:phyx-mini-0307:phyxminiproblems-problemphyxmini0307-answerchoice}
  \lean{PhyXMiniProblems.ProblemPhyXMini0307.AnswerChoice}
  Labels of the four answer choices printed in the source
\end{definition}
\begin{proof}
  This is the declared model definition of Answer Choice.
\end{proof}
\begin{definition}[displayed Reflection Percentage]
  \label{def:physics:phyx-mini-0307:phyxminiproblems-problemphyxmini0307-displayedreflectionpercentage}
  \lean{PhyXMiniProblems.ProblemPhyXMini0307.displayedReflectionPercentage}
  \uses{def:physics:phyx-mini-0307:phyxminiproblems-problemphyxmini0307-answerchoice}
  Percentage printed beside each multiple-choice label
\end{definition}
\begin{proof}
  This is the declared model definition of displayed Reflection Percentage.
\end{proof}
\begin{definition}[reflection Coefficient Percentage]
  \label{def:physics:phyx-mini-0307:phyxminiproblems-problemphyxmini0307-reflectioncoefficientpercentage}
  \lean{PhyXMiniProblems.ProblemPhyXMini0307.reflectionCoefficientPercentage}
  \uses{def:physics:phyx-mini-0307:phyxminiproblems-problemphyxmini0307-boundaryreflectionsetup}
  Convert a dimensionless power reflection coefficient to percent
\end{definition}
\begin{proof}
  This is the declared model definition of reflection Coefficient Percentage.
\end{proof}
\begin{definition}[Matches Answer Choice]
  \label{def:physics:phyx-mini-0307:phyxminiproblems-problemphyxmini0307-matchesanswerchoice}
  \lean{PhyXMiniProblems.ProblemPhyXMini0307.MatchesAnswerChoice}
  \uses{def:physics:phyx-mini-0307:phyxminiproblems-problemphyxmini0307-boundaryreflectionsetup, def:physics:phyx-mini-0307:phyxminiproblems-problemphyxmini0307-answerchoice, def:physics:phyx-mini-0307:phyxminiproblems-problemphyxmini0307-displayedreflectionpercentage, def:physics:phyx-mini-0307:phyxminiproblems-problemphyxmini0307-reflectioncoefficientpercentage}
  The modeled percentage agrees with one of the displayed choices
\end{definition}
\begin{proof}
  This is the declared model definition of Matches Answer Choice.
\end{proof}
\begin{definition}[recorded Dataset Answer]
  \label{def:physics:phyx-mini-0307:phyxminiproblems-problemphyxmini0307-recordeddatasetanswer}
  \lean{PhyXMiniProblems.ProblemPhyXMini0307.recordedDatasetAnswer}
  \uses{def:physics:phyx-mini-0307:phyxminiproblems-problemphyxmini0307-answerchoice}
  Answer label recorded by the source dataset; this is not a premise
\end{definition}
\begin{proof}
  This is the declared model definition of recorded Dataset Answer.
\end{proof}
\begin{lemma}[reflected to incident amplitude ratio]
  \label{lem:physics:phyx-mini-0307:phyxminiproblems-problemphyxmini0307-reflected-to-incident-amplitude-ratio}
  \lean{PhyXMiniProblems.ProblemPhyXMini0307.reflected_to_incident_amplitude_ratio}
  \uses{def:physics:phyx-mini-0307:phyxminiproblems-problemphyxmini0307-boundaryreflectionsetup, def:physics:phyx-mini-0307:phyxminiproblems-problemphyxmini0307-incidentamplitudesi, def:physics:phyx-mini-0307:phyxminiproblems-problemphyxmini0307-reflectedamplitudesi, def:physics:phyx-mini-0307:phyxminiproblems-problemphyxmini0307-matchesproblemandfiguredata, def:physics:phyx-mini-0307:phyxminiproblems-problemphyxmini0307-hasphysicalwaveparameters, def:physics:phyx-mini-0307:phyxminiproblems-problemphyxmini0307-satisfiesstandingwaveinterferencelaw}
  An SWR of 3/2, together with SWR = (A+B)/(A-B), yields the dimensionless amplitude ratio B/A = 1/5. This is an intermediate derived result, not a field of any assumption structure
\end{lemma}
\begin{proof}
  The conclusion follows from the governing relations and figure readouts named in the statement of reflected to incident amplitude ratio.
\end{proof}
\begin{lemma}[reflection Coefficient eq one twenty fifth]
  \label{lem:physics:phyx-mini-0307:phyxminiproblems-problemphyxmini0307-reflectioncoefficient-eq-one-twenty-fifth}
  \lean{PhyXMiniProblems.ProblemPhyXMini0307.reflectionCoefficient_eq_one_twenty_fifth}
  \uses{def:physics:phyx-mini-0307:phyxminiproblems-problemphyxmini0307-boundaryreflectionsetup, def:physics:phyx-mini-0307:phyxminiproblems-problemphyxmini0307-matchesproblemandfiguredata, def:physics:phyx-mini-0307:phyxminiproblems-problemphyxmini0307-hasphysicalwaveparameters, def:physics:phyx-mini-0307:phyxminiproblems-problemphyxmini0307-satisfiesstandingwaveinterferencelaw, def:physics:phyx-mini-0307:phyxminiproblems-problemphyxmini0307-satisfiescommonmediumpowerreflectionlaw}
  The common-medium squared-amplitude power law turns the derived amplitude ratio into the exact power reflection coefficient R = 1/25
\end{lemma}
\begin{proof}
  The conclusion follows from the governing relations and figure readouts named in the statement of reflection Coefficient eq one twenty fifth.
\end{proof}
% --- Archon declaration topology end ---
% --- Archon physics formalization source end ---
